\section{Reglas de Negocio}

Las reglas de negocio establecen las restricciones y directrices operativas que el sistema debe validar para garantizar que los procesos administrativos se realicen bajo los principios de legalidad, transparencia y seguridad institucional de la Defensoría.

% --- REGLAS DE ACCESO Y SEGURIDAD ---
\begin{table}[h!]
\centering
\scriptsize
\caption{Reglas de Negocio: Acceso, Seguridad y Auditoría.}
\label{tab:rn_seguridad}
\begin{tabular}{@{} l p{12cm} @{}}
\toprule
\rowcolor[HTML]{2D3E50} 
\textbf{\color[HTML]{FFFFFF} ID} & \textbf{\color[HTML]{FFFFFF} Regla de Negocio} \\ \midrule
RN-01 & Solo miembros activos de la comunidad politécnica con boleta o número de empleado pueden registrar quejas. \\
\rowcolor[HTML]{F2F4F7} 
RN-02 & El acceso administrativo es restringido por roles; un actor no puede ejecutar funciones de un nivel superior (RBAC). \\
RN-03 & La asignación, alta o modificación de permisos de usuario debe registrarse en un log de auditoría inmutable. \\
\rowcolor[HTML]{F2F4F7} 
RN-04 & Los datos personales del quejoso deben permanecer cifrados en la base de datos y solo ser visibles por personal autorizado. \\
RN-05 & Ningún usuario administrativo puede eliminar registros de auditoría o historiales de transacciones del sistema. \\
\rowcolor[HTML]{F2F4F7} 
RN-06 & El sistema debe cerrar la sesión automáticamente tras 15 minutos de inactividad para prevenir accesos no autorizados. \\ \bottomrule
\end{tabular}
\end{table}

% --- REGLAS DE COMPETENCIA Y REGISTRO ---
\begin{table}[h!]
\centering
\scriptsize
\caption{Reglas de Negocio: Registro y Determinación de Competencia.}
\label{tab:rn_competencia}
\begin{tabular}{@{} l p{12cm} @{}}
\toprule
\rowcolor[HTML]{2D3E50} 
\textbf{\color[HTML]{FFFFFF} ID} & \textbf{\color[HTML]{FFFFFF} Regla de Negocio} \\ \midrule
RN-07 & Toda queja debe incluir obligatoriamente una narrativa de hechos, autoridad imputada y unidad académica. \\
\rowcolor[HTML]{F2F4F7} 
RN-08 & Si el quejoso es menor de edad, el sistema exige el registro de un tutor con identificación oficial válida. \\
RN-09 & No se aceptarán quejas anónimas; la identidad del promovente debe ser validada contra bases de datos del IPN. \\
\rowcolor[HTML]{F2F4F7} 
RN-10 & Las quejas presentadas después de 90 días naturales de ocurrido el hecho se clasificarán como extemporáneas. \\
RN-11 & El sistema no admitirá quejas de carácter laboral o sindical según lo estipulado en el Artículo 13. \\
\rowcolor[HTML]{F2F4F7} 
RN-12 & El motor de IA solo proporciona una sugerencia de competencia; la decisión final es siempre humana. \\
RN-13 & La generación del folio único es automática e irreversible; no puede existir duplicidad de folios en el sistema. \\ \bottomrule
\end{tabular}
\end{table}

% --- REGLAS DE PROCESO Y SUSTANCIACIÓN ---
\begin{table}[h!]
\centering
\scriptsize
\caption{Reglas de Negocio: Proceso, Sustanciación e Investigación.}
\label{tab:rn_proceso}
\begin{tabular}{@{} l p{12cm} @{}}
\toprule
\rowcolor[HTML]{2D3E50} 
\textbf{\color[HTML]{FFFFFF} ID} & \textbf{\color[HTML]{FFFFFF} Regla de Negocio} \\ \midrule
RN-14 & El estatus "En Análisis" se activa únicamente cuando el Recepcionista turna el expediente a la Subdefensoría. \\
\rowcolor[HTML]{F2F4F7} 
RN-15 & El sistema debe bloquear el avance del folio si no se ha cargado el Acuerdo de Admisión correspondiente. \\
RN-16 & Todo expediente radicado debe tener asignado obligatoriamente a un abogado analista responsable. \\
\rowcolor[HTML]{F2F4F7} 
RN-17 & Los oficios enviados a Unidades Académicas deben generar automáticamente un recordatorio de respuesta. \\
RN-18 & El plazo de respuesta otorgado a la autoridad debe ser de 2 a 15 días hábiles según la normativa. \\
\rowcolor[HTML]{F2F4F7} 
RN-19 & No se puede pasar a la fase de resolución si existen actos de investigación con plazos de respuesta vigentes. \\
RN-20 & Las notas internas del Analista no son descargables ni visibles para el quejoso en el portal público. \\
\rowcolor[HTML]{F2F4F7} 
RN-21 & El estatus "En Investigación" es obligatorio antes de emitir cualquier propuesta de conclusión de fondo. \\ \bottomrule
\end{tabular}
\end{table}

% --- REGLAS DE CIERRE Y FIRMA DIGITAL ---
\begin{table}[h!]
\centering
\scriptsize
\caption{Reglas de Negocio: Conclusión, Firma y Archivo.}
\label{tab:rn_cierre}
\begin{tabular}{@{} l p{9.5cm} l @{}}
\toprule
\rowcolor[HTML]{2D3E50} 
\textbf{\color[HTML]{FFFFFF} ID} & \textbf{\color[HTML]{FFFFFF} Regla de Negocio} \\ \midrule
RN-22 & El Acuerdo de Conclusión debe ser visado por el Subdefensor antes de enviarse a firma del Titular. \\
\rowcolor[HTML]{F2F4F7} 
RN-23 & La notificación al quejoso de la propuesta de conclusión es requisito previo para el cierre del folio. \\
RN-24 & La Firma Digital Avanzada del Defensor es el único medio para validar legalmente el cierre del caso. \\
\rowcolor[HTML]{F2F4F7} 
RN-25 & El sistema impedirá la modificación de cualquier documento una vez que haya sido firmado digitalmente. \\
RN-26 & El estatus "Finalizado" solo se activa mediante la aplicación exitosa de la firma electrónica del Defensor. \\
\rowcolor[HTML]{F2F4F7} 
RN-27 & Todo expediente finalizado debe ser transferido al archivo histórico en un periodo no mayor a 24 horas. \\
RN-28 & El envío a archivo histórico es irreversible y bloquea cualquier edición posterior de los datos del caso. \\ \bottomrule
\end{tabular}
\end{table}

% --- REGLAS DE GESTIÓN Y ADMINISTRACIÓN ---
\begin{table}[h!]
\centering
\scriptsize
\caption{Reglas de Negocio: Estadísticas y Administración.}
\label{tab:rn_admin}
\begin{tabular}{@{} l p{12cm} @{}}
\toprule
\rowcolor[HTML]{2D3E50} 
\textbf{\color[HTML]{FFFFFF} ID} & \textbf{\color[HTML]{FFFFFF} Regla de Negocio} \\ \midrule
RN-29 & Solo el rol de Defensor (Titular) tiene facultad para visualizar tableros de control global y estadísticas. \\
\rowcolor[HTML]{F2F4F7} 
RN-30 & Las métricas de reincidencia deben calcularse automáticamente basándose en los datos históricos del repositorio. \\
RN-31 & La configuración de plantillas de documentos oficiales es de acceso restringido al perfil Administrador. \\
\rowcolor[HTML]{F2F4F7} 
RN-32 & El sistema debe generar alertas preventivas ante el rezago administrativo (expedientes sin movimiento en 10 días). \\
RN-33 & Las copias de seguridad de los expedientes deben almacenarse en un repositorio inmutable y externo al sistema. \\
\rowcolor[HTML]{F2F4F7} 
RN-34 & Toda actualización de estado del folio debe generar una notificación automática al correo del quejoso. \\
RN-35 & El catálogo de Unidades Académicas y dependencias debe estar sincronizado con la estructura oficial del IPN. \\
\rowcolor[HTML]{F2F4F7} 
RN-36 & No se pueden emitir oficios si la autoridad responsable no se encuentra debidamente registrada en el sistema. \\
RN-37 & Los acuerdos de prevención deben otorgar al quejoso un plazo improrrogable de 5 días para subsanar omisiones. \\
\rowcolor[HTML]{F2F4F7} 
RN-38 & El rechazo de una queja por incompetencia debe estar fundado obligatoriamente en los criterios del Artículo 13. \\
RN-39 & Las propuestas de conciliación requieren la firma de aceptación del quejoso cargada en el sistema para proceder. \\
\rowcolor[HTML]{F2F4F7} 
RN-40 & El sistema debe garantizar que el folio digital se mantenga accesible para consulta del quejoso por un año mínimo. \\ \bottomrule
\end{tabular}
\end{table}